\documentclass[11pt]{article}
\usepackage[T1]{fontenc}
\usepackage{textcomp}
\usepackage[margin=0.75in]{geometry}
\usepackage{amsmath,amssymb,amsthm}
\usepackage{array,booktabs}
\usepackage{hyperref}
\setlength{\parindent}{0pt}
\newtheorem{theorem}{Theorem}
\theoremstyle{definition}
\newtheorem{definition}{Definition}
\title{Flow Proof Artifact: Finset-union-self-card-bound-derived}
\author{Generated by \texttt{flow doc proof}}
\date{Flow Proof Book}
\begin{document}
\maketitle
Union with self does not increase cardinality.\\[0.5em]
\textbf{Source.} Finset union self card bound\\[0.5em]
\bigskip
\noindent{\large \textbf{Derived fact 1.} \textit{Union with self does not increase cardinality} \hfill Finset \cdot union \cdot self card bound derived}\par
\smallskip\noindent\textit{Coordinate: Finset \cdot union \cdot self card bound derived}\par
\smallskip\noindent\textit{Source: Finset union self card bound}\par
\smallskip\noindent\textit{Built on: \hyperref[thm:Finset-union-self_card_bound_derived]{Derived fact 1}}\par
\medskip
\noindent\textbf{Goal.} Union with self does not increase cardinality\par
\noindent$\displaystyle \forall s \in \guillemotleft{}Finset\guillemotright{} \guillemotleft{}Nat\guillemotright{}\quad «Finset» «card»(«Finset» «union»(s, s)) \le «Finset» «card»(s)$\par
\medskip
\noindent
\begin{tabular}{@{} >{\raggedright\arraybackslash}p{0.06\textwidth} >{\raggedright\arraybackslash}p{0.47\textwidth} >{\raggedleft\arraybackslash}p{0.41\textwidth} @{}}
\textbf{\#} & \textbf{Proof} & \textbf{Mathematics} \\
\midrule
\textbf{1.} & We prove that self card bound derived for union on Finset. & \\[0.45em]
\textbf{2.} & We invoke the derived fact: \guillemotleft{}Finset\guillemotright{} \guillemotleft{}card\guillemotright{} (instantiated for \guillemotleft{}Finset\guillemotright{} \guillemotleft{}union\guillemotright{}(s, s)) <= \guillemotleft{}Finset\guillemotright{} \guillemotleft{}card\guillemotright{}(s). & \\[0.45em]
\textbf{3.} & From step 2, this implies \guillemotleft{}Finset\guillemotright{} \guillemotleft{}card\guillemotright{}(\guillemotleft{}Finset\guillemotright{} \guillemotleft{}union\guillemotright{}(s, s)) is at most \guillemotleft{}Finset\guillemotright{} \guillemotleft{}card\guillemotright{}(s). Hence proven. & $\displaystyle «Finset» «card»(«Finset» «union»(s, s)) \le «Finset» «card»(s)$ \\[0.45em]
\bottomrule
\end{tabular}
\label{thm:Finset-union-self_card_bound_derived}
\par\medskip\hrule
\end{document}
